\documentclass[aspectratio=169]{beamer}
\input{header}
\coursecode{JKSSB}

\title{Web Browsing, Search, Upload \& Download}
\subtitle{Lesson 36 --- Search, Storage and Transfer}
\author{JKSSB Computer Awareness}
\date{\today}

\begin{document}
\frame{\titlepage}

\begin{frame}{What a Web Browser Does}
  \begin{itemize}
    \item A \key{web browser} is the software used to open and view web pages,
      such as Google Chrome, Mozilla Firefox, and Microsoft Edge.
    \item The browser sends a request for a page and displays the returned
      text, images, and links.
    \item The page address typed in the \key{address bar} is called the
      \key{URL} --- Uniform Resource Locator.
    \item Each open page lives in its own \key{tab}, so several pages can
      stay open at once.
  \end{itemize}
  \begin{block}{The one idea}
    The browser fetches pages; the search engine finds pages; tabs keep
    several pages open together.
  \end{block}
\end{frame}

\begin{frame}{Tabs, Bookmarks, and History}
  \begin{itemize}
    \item \key{Tabs} let many pages stay open in one window; a new tab opens
      with \key{Ctrl+T}.
    \item A \key{bookmark} is a saved shortcut to a page so it can be opened
      again without typing the address.
    \item \key{Browsing history} is the list of pages visited earlier,
      stored by date.
    \item History answers ``where was I''; a bookmark answers ``where do I
      want to return''.
  \end{itemize}
  \begin{exampleblock}{Office desktop example}
    While preparing a report, the clerk keeps the office portal in one tab
    and the search results in another, and bookmarks the portal for tomorrow.
  \end{exampleblock}
\end{frame}

\begin{frame}{How a Search Engine Works: Crawling}
  \begin{itemize}
    \item A \key{search engine} (for example Google, Bing) finds pages that
      match the typed words.
    \item Step 1 --- \key{crawling}: automatic programs called
      \key{crawlers} or spiders visit pages and follow their links.
    \item Crawling discovers new and changed pages across the web.
    \item Without crawling, the engine would never learn that a new page
      exists.
  \end{itemize}
  \begin{block}{Pipeline position}
    Crawling is the collection step: it gathers pages before anything else
    can happen.
  \end{block}
\end{frame}

\begin{frame}{Indexing and Ranking}
  \begin{itemize}
    \item Step 2 --- \key{indexing}: the crawled pages are stored in a large
      \key{index}, like the index at the back of a book.
    \item Step 3 --- \key{ranking}: when words are typed, the engine picks
      matching pages from the index and orders them by relevance.
    \item The ordered list of matches is the \key{search results page}.
    \item \muted{Order to remember: crawl, then index, then rank.}
  \end{itemize}
  \begin{alertblock}{Trap alert}
    The index is built \key{before} the search; ranking happens only after
    the words are typed.
  \end{alertblock}
\end{frame}

\begin{frame}{Sponsored Results vs Organic Results}
  \begin{itemize}
    \item \key{Organic results} are the normal ranked matches chosen by
      relevance.
    \item \key{Sponsored results} (labelled Ad / Sponsored) are paid
      placements shown because an advertiser paid, not because they ranked
      highest.
    \item A sponsored result always carries a visible \key{Ad} or
      \key{Sponsored} label.
    \item Being listed first on the page does \key{not} always mean being
      the most relevant match.
  \end{itemize}
  \begin{block}{Keep the pair straight}
    Organic = ranked by relevance; sponsored = placed by payment.
  \end{block}
\end{frame}

\begin{frame}{Cookies: What They Remember}
  \begin{itemize}
    \item A \key{cookie} is a small file a website stores on the computer to
      remember that visitor.
    \item Cookies keep a user logged in, remember items in a cart, and store
      preferences such as language.
    \item They are set \key{by websites} and read back on the next visit.
    \item Deleting cookies signs the user out of sites and clears saved
      preferences.
  \end{itemize}
  \begin{block}{One-line definition}
    A cookie remembers \key{who the visitor is} for the website.
  \end{block}
\end{frame}

\begin{frame}{Cache: What It Stores}
  \begin{itemize}
    \item The browser \key{cache} stores \key{copies of page parts} ---
      images, styles, and scripts --- on the computer.
    \item On the next visit the saved copies load from the disk instead of
      being fetched again, so the page opens \key{faster}.
    \item Cache is managed \key{by the browser}, not by the website.
    \item Clearing the cache does not sign the user out; it only makes the
      next visit slightly slower.
  \end{itemize}
  \begin{block}{One-line definition}
    The cache remembers \key{page parts} so repeat visits load faster.
  \end{block}
\end{frame}

\begin{frame}{Cookies vs Cache vs History}
  \begin{center}
    \begin{tabular}{lll}
      \toprule
      \textbf{Item} & \textbf{What it stores} & \textbf{Who stores it} \\
      \midrule
      Cookie & Visitor data (login, cart, choices) & Website \\
      Cache & Copies of page files (images, styles) & Browser \\
      History & List of visited page addresses & Browser \\
      \bottomrule
    \end{tabular}
  \end{center}
  \vspace{0.3cm}
  \begin{alertblock}{Trap alert}
    Cookie = who you are; cache = saved page parts; history = where you went.
    Do not swap the three.
  \end{alertblock}
\end{frame}

\begin{frame}{Downloading}
  \begin{itemize}
    \item \key{Downloading} means copying a file \key{from the internet to
      the computer}.
    \item Examples: saving a form PDF from the office portal, saving a
      photo from a web page.
    \item A download moves data \key{down} to the user's machine and is
      measured by \key{download speed}.
    \item The saved file then opens without needing the internet.
  \end{itemize}
  \begin{exampleblock}{Office desktop example}
    The clerk downloads the leave-application PDF from the portal to the
    desktop, then fills it offline in a reader.
  \end{exampleblock}
\end{frame}

\begin{frame}{Uploading}
  \begin{itemize}
    \item \key{Uploading} means sending a file \key{from the computer to the
      internet}.
    \item Examples: attaching a document to an email, submitting a scanned
      form on a portal.
    \item An upload moves data \key{up} from the user's machine to a server.
    \item Sending the file needs the internet connection at sending time.
  \end{itemize}
  \begin{alertblock}{Trap alert}
    Download = internet to computer; upload = computer to internet.
    The direction is always judged from the user's machine.
  \end{alertblock}
\end{frame}

\begin{frame}{Download Speed vs Upload Speed}
  \begin{center}
    \begin{tabular}{lll}
      \toprule
      \textbf{Point} & \textbf{Download} & \textbf{Upload} \\
      \midrule
      Direction & Server to computer & Computer to server \\
      Use & Opening pages, saving files & Sending mail attachments, submitting forms \\
      Typical speed & Faster & Slower \\
      \bottomrule
    \end{tabular}
  \end{center}
  \vspace{0.25cm}
  \muted{Home and office plans usually give a faster download than upload.}
\end{frame}

\begin{frame}{File Conversion}
  \begin{itemize}
    \item \key{File conversion} means changing a file from one format to
      another, keeping the same content.
    \item Common office need: converting a Word document (\key{.docx}) to
      \key{PDF} before submitting it on a portal.
    \item Conversion changes the \key{format}, not the transfer direction:
      a converted file must still be uploaded or downloaded separately.
    \item The file \key{extension} (such as .pdf, .docx, .jpg) names the
      format the file is stored in.
  \end{itemize}
  \begin{block}{Extension vs format}
    The extension is the name tag (.png); the format is the storage method
    (PNG). Do not conflate them.
  \end{block}
\end{frame}

\begin{frame}{Safe and Correct Practice}
  \begin{itemize}
    \item Download files only from trusted sites; a free download can carry
      a \key{virus}.
    \item Check the file extension before opening: a document should not
      arrive as an unknown program file.
    \item Use \key{Ctrl+D} to bookmark the current page and \key{Ctrl+H} to
      open browsing history.
    \item Clear history to hide visited addresses; clear cookies to sign
      out of sites.
  \end{itemize}
\end{frame}

\begin{frame}{Trap Alert: Facts to Keep Straight}
  \begin{itemize}
    \item Search order: \key{crawl}, then \key{index}, then \key{rank}.
    \item Sponsored = paid placement with an \key{Ad} label; organic =
      ranked by relevance.
    \item Cookie (website, visitor data) vs cache (browser, page copies)
      vs history (browser, visited list).
    \item Download = to the computer; upload = from the computer to the
      server.
    \item Conversion changes the \key{format}; transfer changes the
      \key{location}.
  \end{itemize}
\end{frame}

\begin{frame}{Lesson 36: Recall}
  \begin{itemize}
    \item Browsers fetch pages; tabs hold several pages; bookmarks save
      shortcuts; history lists visited pages.
    \item Search pipeline: crawlers collect pages, the index stores them,
      ranking orders the matches; sponsored results are paid Ads.
    \item Cookies remember the visitor (website); cache keeps page copies
      (browser); history records visited addresses (browser).
    \item Downloading copies from the internet to the computer; uploading
      sends from the computer to the internet.
    \item File conversion changes the format (for example .docx to PDF);
      the extension names that format.
  \end{itemize}
\end{frame}
\end{document}
